% !TeX root = ../main.tex
% ============================================================
%  sections/02_literature_review.tex
% ============================================================

\section{Literature Review}
\label{sec:literature}

This section surveys three streams of research that converge in our
contribution: (i) bilevel optimisation and its computational
treatment, (ii) closed-loop and sustainable supply chain network
design, and (iii) carbon pricing in operations management and
environmental economics. We close the section by articulating the
research gap that motivates the present paper.

\subsection{Bilevel optimisation}
\label{sec:lit:bilevel}

Bilevel programming, formalised by \citet{bracken1973} in the
context of Stackelberg games \citep{vonStackelberg1934}, has matured
into a substantial subfield of mathematical programming. Foundational
treatments are provided in the monograph of \citet{dempe2002} and the
survey of \citet{colson2007}. Even in its simplest, fully linear
form, bilevel optimisation is NP-hard \citep{jeroslow1985, hansen1992},
which has motivated a steady stream of methodological advances.

\subsubsection{Exact reformulations}
\label{sec:lit:exact}

The most common exact route
replaces the lower-level problem by its KKT optimality conditions and
linearises the resulting complementarity constraints
\citep{fortuny1981}. While this approach is conceptually simple, the
disjunctive structure of complementarity rapidly degrades scalability,
particularly when the follower contains binary variables. Tailored
branch-and-bound and branch-and-cut methods that exploit the
hierarchical structure have been developed by \citet{moore1990} and,
in the mixed-integer case, by \citet{denegre2011}.

\subsubsection{Bounding and decomposition}
\label{sec:lit:bounding}

Penalty-based bounds for
bilevel linear programs were proposed by \citet{edmunds1991}.
Subsequent work has produced general-purpose exact solvers for mixed-
integer bilevel linear programs (MIBLPs): \citet{fischetti2017}
develop a branch-and-cut algorithm based on intersection and no-good
cuts that solves hundreds of previously open benchmark instances;
\citet{lozanosmith2017} pursue an alternative, value-function-based
reformulation that avoids an explicit single-level recast of the
follower; and \citet{tahernejad2020} describe a generalised
branch-and-cut framework, implemented in the open-source solver
\textsc{MibS}, that unifies several of these cutting strategies.
\citet{kleinert2021} survey this rapidly developing area,
highlighting the role of no-good cuts, disjunctive cuts, and
intersection cuts in mixed-integer bilevel programming. Our \DBP is
in a similar spirit --- it bounds the leader's optimal value from
below and above and tightens both bounds iteratively --- but it
exploits the specific structure of a \emph{discrete, finite} leader
domain (a tax on a grid) coupled with a parametric mixed-integer
follower, rather than relying on generic disjunctive cuts over a
continuous leader space.

\subsubsection{Bilevel models for network design and pricing}
\label{sec:lit:bilevel_network}

Bilevel optimisation has a long history in network design and
pricing problems, where a leader sets tolls, tariffs, or capacities
and a follower (or population of followers) reacts through flow or
allocation decisions; \citet{brotcorne2001} is a classical example
for toll optimisation on a multicommodity transportation network.
Closer to our setting, \citet{yueyou2017} formulate a Stackelberg
game between a supply-chain leader and operator as a mixed-integer
bilevel program in which the follower retains discrete facility,
technology-selection, and shutdown decisions, solved through a
reformulation-and-decomposition algorithm. Their framework establishes
that a discrete-follower bilevel \SCND model is computationally
tractable, but it is not posed as an environmental-regulation problem
and the leader does not face a finite policy grid; we build on this
precedent in \Cref{sec:lit:gap}.

\subsection{Closed-loop supply chain network design}
\label{sec:lit:scnd}

The \SCND literature is extensive; comprehensive reviews are given
by \citet{melo2009}, \citet{farahani2014}, and, with explicit
attention to sustainability, by \citet{eskandarpour2015} and
\citet{govindan2015}.

\subsubsection{Closed-loop models}
\label{sec:lit:closedloop_models}

\CLSC models jointly optimise
forward and reverse logistics, capturing collection,
remanufacturing, and reuse alongside production and distribution
decisions. \citet{savaskan2004} provide the canonical analytical
treatment of channel choice for collecting used products for
remanufacturing; \citet{fleischmann2001} propose an early integrated
network-design formulation. Subsequent work extends these models along
several dimensions: to multi-period, multi-product, and multi-modal
settings \citep{pishvaee2010, ozceylan2014}; to the joint design of
forward and reverse flows under demand uncertainty
\citep{cardoso2013}; to multi-period capacitated location with
delayed demand satisfaction \citep{correia2014} and $p$-robust
location under disruption \citep{snyder2006}; and to a triple-bottom-
line treatment of cost, emissions, and social impact solved via
hybrid metaheuristics \citep{devika2014}. Robust and stochastic
extensions addressing demand and return-rate uncertainty are
presented by \citet{pishvaee2011} and \citet{keyvanshokooh2016}. Our
follower formulation builds on this stream and extends it through
intermodal hub routing in the spirit of strategic hub-network design
\citep{yildiz2021}, technology-specific production and
remanufacturing options, and explicit treatment of storage-area
capacity.

\subsubsection{Carbon aware SCND}
\label{sec:lit:carbon_aware}

A growing body of work embeds
emission considerations into network design through emission caps,
multi-objective formulations trading off cost and emissions, or
explicit tax payments. \citet{chaabane2012} and
\citet{fahimnia2013} analyse the impact of carbon policies on
distribution networks; \citet{mota2018} integrate sustainability
indicators into a multi-product \CLSC model with technology
selection; and \citet{zhen2019} pose a bi-objective cost-versus-
CO\textsubscript{2} closed-loop network design problem under demand
uncertainty. In every one of these studies, however, the carbon
price or emission cap is a fixed parameter set by the modeller, not a
decision optimised by an explicit regulator.

\subsection{Carbon pricing in operations management and
environmental economics}
\label{sec:lit:carbon}

The choice between price-based instruments (a tax) and quantity-based
instruments (a cap or tradable permits) for controlling an
externality is a classical question in environmental economics.
\citet{weitzman1974} shows that the relative efficiency of prices
versus quantities depends on the curvature of marginal costs and
benefits, and \citet{newellpizer2003} extend the comparison to stock
externalities such as greenhouse gases under uncertainty, a setting
in which price instruments are generally favoured. \citet{aldy2010}
synthesise the resulting policy-design literature, covering
instrument choice, stringency, and dynamic efficiency for climate
mitigation. These results motivate our leader's choice of a tax
\emph{level} on a fixed instrument (rather than a quantity cap) as
the policy variable, and they underlie the over-taxation pathology we
document in \Cref{sec:exp:tradeoff}: a price instrument calibrated
purely to a quantity (emissions) target is, in Weitzman's terms,
mismatched to a follower whose marginal abatement cost is itself
discontinuous and eventually near-vertical.

Within operations management, \citet{tang2015} and
\citet{benjaafar2013} compare carbon tax, cap-and-trade, and
cap-and-offset schemes, characterising their respective effects on
production, inventory, and supply-chain decisions. \citet{xu2016}
study a two-echelon supply chain operating under a carbon tax, while
\citet{ji2017} examine a single-manufacturer \SCND problem under
cap-and-trade, and \citet{zakeri2015} investigate the joint design of
forward and reverse flows under emissions regulation. None of these
papers, however, treats the tax level itself as a decision made by an
optimising regulator who anticipates the supply chain's response.

\subsubsection{Bilevel models with environmental regulation}
\label{sec:lit:bilevel_env}

A small but growing literature applies bilevel optimisation to
environmental policy design. \citet{zhang2021} consider a bilevel
model for carbon-quota allocation but restrict the follower to a
linear program with no integer variables, foreclosing facility-level
decisions. \citet{he2019} develop a bilevel green supply-chain model
focused on pricing rather than network design.
\citet{cao2020} address tax differentiation across regions but treat
the supply-chain problem as a continuous flow problem.

\subsection{Research gap and positioning}
\label{sec:lit:gap}

Four observations emerge from the surveyed literature.
\emph{First}, the vast majority of carbon-aware \SCND studies
(\Cref{sec:lit:scnd}) model the tax or cap as a fixed parameter and thus
cannot inform the regulator's choice of an effective tax level.
\emph{Second}, the existing bilevel formulations of environmental
regulation \citep{zhang2021, he2019, cao2020} simplify the follower
problem to a continuous, single-period setting, omitting integer
facility, technology, and capacity decisions that are central to
long-run emissions; conversely, the bilevel \SCND model that
\emph{does} retain a discrete follower \citep{yueyou2017} is not
posed as an environmental-policy problem and has no regulator
objective over emissions. \emph{Third}, environmental economics
clarifies \emph{why} the leader's instrument matters
\citep{weitzman1974, newellpizer2003, aldy2010} but stops at the
single-firm, reduced-form level and does not model a multi-echelon
network operator as the follower. \emph{Fourth}, even when a bilevel
mixed-integer model is proposed elsewhere in the optimisation
literature, algorithmic treatment is often limited to KKT
reformulations or generic disjunctive cuts
\citep{fortuny1981,fischetti2017,kleinert2021} whose scalability or
applicability has not been demonstrated for a discrete leader domain
coupled with a multi-period closed-loop follower.

To the best of our knowledge, no prior work simultaneously
(i) formulates a \emph{bilevel} carbon-tax problem with a
\emph{multi-echelon, multi-period closed-loop} \SCND follower
incorporating integer facility, technology, capacity, and
transportation decisions; and (ii) develops a convergent
\emph{deterministic bounding procedure} adapted to the discrete
leader domain, together with a lexicographic tie-break that exposes
--- rather than conceals --- the over-taxation pathology implied by
\citet{weitzman1974}'s price-versus-quantity argument. The present
paper fills this gap.
